% Cote 119: offprints, reports, courses, CV and letters (typed and printed part): archive cote 119, site page 258 (PDF page 259).
% typescript. Draft compared: none.
% AI-transcribed by Claude Opus 5.5 (claude-opus-5-5, adaptive thinking, effort high), Message Batches API, batch r119_full (msgbatch_01DwZNm5Y7YsKBWyAN4sFs8w), 2026-11-13; not yet checked by eye.
\scanpage{119/258}
\mspage{5}

gnement de très grande qualité, y compris à un niveau élevé, dans une Univer-%
sité qui manque cruellement de telles forces, ils se trouvent écartés de telles
fonctions auxquelles tout les prédispose. D'autres, parfaitement capables de
faire des enseignants honorables, mais n'ayant pas les dons voulus pour faire
des travaux originaux, feront appel à des ``patrons'' réputés particulièrement
peu exigeants pour faire passer une thèse, et avec l'appui de ceux-ci se
verront nommés maîtres de conférences,alors que tel de leurs camarades, plus
qualifié mais plus scrupuleux, quittera le CNRS pour devenir simple assistant,
ou cherchera un poste dans l'enseignement secondaire. Par la suite, de tels
maîtres de conférences, qualifiés certes pour enseigner à l'Université, mais
peu ou point pour juger de questions concernant la formation ou le recrutement
de mathématiciens chercheurs, auront à ce sujet voix au chapitre au même
titre que ceux de leurs collègues ayant fait des contributions appréciables
à la science, et seront habilités pour faire passer des thèses (soi-disant de
recherche) tout comme ceux-ci. On voit que faute d'une définition des compé-%
tences qui corresponde vraiment à la réalité, des hommes ayant une qualifica-%
tion professionnelle sérieuse, qui pourrait sans usurpation aucune de leur
part en faire des membres utiles et à part entière de la communauté univer-%
sitaire, peuvent se trouver amenés, par la logique même du système dans
lequel ils sont placés et dont ils tirent parti à leur façon, dans une
situation où ils risquent de fausser la vie scientifique proprement dite
dans leur département, voire,dans des cas extrêmes, de l'étouffer.

3. Compte tenu de tout ce qui précède, il apparaît donc urgent de redéfinir,
les rôles de maître-enseignant et de maître-chercheur, pour les dissocier
tout en tenant compte en même temps de leurs interrelations étroites.

Le \uline{maître-chercheur} est un homme qui s'intéresse en premier lieu
aux problèmes que lui pose la science, plus qu'aux hommes qui l'apprennent
ou l'utilisent, plus même qu'aux hommes qui la font. L'enseignement souvent
pour lui est surtout un moyen, lui assurant une position sociale qui lui
permette de se consacrer à l'édification de la science indépendamment des
applications pratiques qui peuvent la rendre immédiatement rentable. Il est,

\margin{258}
